%	현대물리실험 보고서
%	실험3
%	202100973 이승엽

%----------------------------------------------------------------------------------------
%	PACKAGES AND DOCUMENT CONFIGURATIONS
%----------------------------------------------------------------------------------------

\documentclass[a4paper, 10pt, nanum]{CSUniSchoolLabReport}
% use UTF8 encoding
\usepackage[utf8]{inputenc}
% use KoTeX package for Korean
\usepackage{kotex}
\usepackage{amsmath}
\usepackage{hyperref}
\usepackage{graphicx}
\usepackage{indentfirst}
\usepackage{setspace}
\usepackage{multirow}
\usepackage{enumitem}
\usepackage{graphicx}
\usepackage{wrapfig}
\usepackage{epstopdf}
\usepackage{caption}
\usepackage{subcaption}

\setlength{\parindent}{0.2in} % 들여쓰기 길이 설정
\setlength{\parskip}{3mm} % 문단 간의 간격 조절
\setstretch{1.5} % 줄간격

\addbibresource{reference.bib} % Bibliography file (located in the same folder as the template)

%----------------------------------------------------------------------------------------
%	REPORT INFORMATION
%----------------------------------------------------------------------------------------

\title{현대물리실험 실험3 보고서 \\ Franck-Hertz experiment with neon} % Report title

\author{\textsc{Department of Physics} 202100973 이승엽}

\date{\today}

%----------------------------------------------------------------------------------------

\begin{document}

\maketitle % Insert the title, author and date using the information specified above

\begin{center}
	\begin{tabular}{l r}
		Date Performed: & April 11, 2023 \\ % Date the experiment was performed
		& April 18, 2023 \\
		Partners: & 202100969 이규리 \\ % Partner names
		& 202100989 한누리 \\
		Instructor: & Professor 이기주 \\ % Instructor/supervisor
		Typesetting: & LaTeX \\
		Datafitting: & Python \\
	\end{tabular}
\end{center}

%----------------------------------------------------------------------------------------
%	ABSTRACT
%----------------------------------------------------------------------------------------

\maketitle
% \begin{abstract}
% 	This report ...
% \end{abstract}

%----------------------------------------------------------------------------------------
%	INTRODUCTION
%----------------------------------------------------------------------------------------

\section{서론}

  To record a Franck-Hertz curve for neon. To measure the discontinuous energy emission of free electrons for inelastic collision. To interpret the measurement results as representing discrete energy absorption by neon atoms. To observe the Ne-spectral lines resulting from the electron-collision excitation of neon atoms. To identify the luminance phenomenon as layers with a high probability of excitation. 

%----------------------------------------------------------------------------------------
%	THEORY
%----------------------------------------------------------------------------------------

\section{이론}

	\begin{figure}[htb!]
		\centering
		\includegraphics[width=5cm]{fig1.png}
		\caption{Top : Simplified term diagram for neon.\\ Bottom : the electron current flowing to the collector as a function of the acceleration voltage in the Franck-Hertz
		experiment with neon}
		\label{fig:1}
	\end{figure}
  J. Frank와 G. Hertz가 1914년 이후 원자의 공명 퍼텐셜을 구하기 위해 실시한 실험으로서, 수은이 네온으로 대체된 것만 빼면 같은 실험으로 당시 보어의 수소 원자 모형이 가진 전자의 Life time 문제를 설명하는 에너지 준위의 양자화를 직접적으로 보여준 실험이다. 이 실험으로 같은 해에 제창된 N.H.D. 보어의 원자 구조론에 대하여 실험적 근거를 밝혀 양자론 발전에 기여하였다. 열전자를 가속하여 전자빔을 만들어 저
  압 기체 속을 지나가게 하여 그 원자와 충돌시키도록 한자. 이 전자흐름에 의한 전류의 세기를 조사하면 전자가 원자와의 충돌에 의해 어떤 에너지의 변화를 받는가를 알 수 있다. 실험결과 , 가속 전압이 점차로 증가하여 어느 한계에 도달하면 전류가 급격히 감소
  하고, 가속 전압이 증가하면 전류는 다시 증가하여 전압이 제2차 한계 값에 이르면 다시 급속히 감소하고, 이후에는 이와 같은 전류의 극대 값이 반복된다. 이와 같이, 네온 원자에 속박된 전자의 에너지 준위가 양자화 된 것을 확인함을 목적으로 한다.

%----------------------------------------------------------------------------------------
%	EXPERIMENTAL METHOD
%----------------------------------------------------------------------------------------

\section{실험 방법}

	ㅇ

%----------------------------------------------------------------------------------------
%	RESULT AND DISCUSSION
%----------------------------------------------------------------------------------------

\section{실험 결과}

	ㅇ


%----------------------------------------------------------------------------------------
%	CONCLUSIONS
%----------------------------------------------------------------------------------------

\section{결론}

	ㅇ

%----------------------------------------------------------------------------------------
%	BIBLIOGRAPHY
%----------------------------------------------------------------------------------------

% \printbibliography % Output the bibliography

%----------------------------------------------------------------------------------------

\end{document}